\section{REFERÊNCIAL TEÓRICO}

Esta seção apresenta conceitos necessários para o desenvolvimento do trabalho, sendo esses: definição de conceitos básicos (linguagens de programação e modelo cliente-servidor) descrição das tecnologias usadas para a construção da aplicação (\en{Java}, \en{Kotlin}, redes \en{Peer-to-Peer}), diferentes tipos de estratégias de sincronização e tópicos essenciais para o seu entendimento tão como comparações com trabalhos relacionados. Segue a descrição geral dos tópicos abordados.

\en{Middleware} forma uma camada entre aplicações e plataformas distribuídas para prover transparência de distribuição, isso é, esconder uma porção dos dados, processamento e controle das aplicações \cite{tanenbaum2007dist}. Desse modo, será apresentado os conceitos necessários para montar essa camada intermediária.

As formas de classificar as estratégias de sincronização podem ser de acordo com a necessidade da aplicação em relação a consistência entre as réplicas \cite{tanenbaum2007dist}. Assim, será destacado definições de conceitos essenciais para a categorização e implementação das estratégias de sincronização de dados. 

A aplicação consiste em um leitor de \en{ePub} com a possibilidade de fazer anotações, a sincronização será feita a partir dos compartilhamento dessas anotações na rede, todos os pares devem ter as anotações dos demais. Desse modo, as tecnologias relevantes são: \en{Java} e \en{Kotlin} para linguagens de programação e \en{Room} para banco de dados local, onde os dados serão armazenados no dispositivo móvel.

\subsection{SISTEMAS DISTRIBUÍDOS E A INTERNET}

\begin{comment}
Adicionar:
  - Relacionar contexto do ISAM com o contexto da INTERNET

Melhorar:
  - Pegar funcionalidades do CORBA quando listar as funcionalidades de middlewares
  - Referência para embasar a afirmação sobre os componentes dos sitemas distribuídos

Estrutura do paragrafo
  - Contextualizar a internet como um conjunto de sistemas distribuídos
  - Papel dos protocolos na internet
  - Usar exemplos de aplicações reais, como editores compartilhados
  - Como peer-to-peer se relaciona com a internet (DNS, IPv4, IPv6)
  - Indicar as relações dos servidores root com a comunicação peer-to-peer
\end{comment}


A \en{internet} ocupa um papel essencial na atualidade, através da computação móvel, aplicações empresariais, redes \en{LAN}, redes sem-fio (\en{WIFI}), e outros exemplos de conectividade que possibilitam a conexão. Essas conexões ocorrem de diversas formas: \en{WiMAX} (interoperabilidade mundial de acesso a micro-ondas), \en{Bluetooth}, quarta geração de redes móveis (4G) e outros exemplos de trammissão de dados pelo meio físico. Para possibilitar essa comunicação há os protocolos de rede como o \en{TCP} e o \en{IP} que, dentre outros, formam o modelo referêncial \en{TCP/IP} que define a pilha de protocolos reponsáveis por ligar os dados vindos pelos meios físicos às aplicações \cite{alberti2016distributedsystemsfutureinternet}. Nota-se, desse modo, como sistemas distintos conseguem cooperar na \en{internet}.

Esse conjunto de sistemas distintos dividem a tarefa de processamento de dados entre si, formando sistemas distribuídos. Um sistema distribuído é um ambiente de computação, onde vários componentes localizados em multiplos dispositivos de computação em uma rede segregam trabalho entre si, de modo que, para o usuário final, pareça ser um sistema único que realiza o trabalho \cite{khole2023compendium}. Os componentes de um sistema distribuídos são \en{softwares} com responsabilidades específicas, como: \en{APIs} que são resposáveis por conectar o cliente com o componente de persistência de dados e aplicar a lógica de negócio; \en{middlewares} responsáveis por afunilar o trafégo de dados, filtrando dados, aplicando restrição de acesso, validações em requisições, etc. 


\subsection{LINGUAGENS DE PROGRAMAÇÃO}

A linguagem nativa de computadores é a binária (uns e zeros), todas as instruções e dados devem ser providas nesse formato para serem processados. Código binário nativo é chamado de linguagem de máquina. Com o passar do tempo foram aparecendo novas formas de construir programas; sucedendo o código de máquina os programas começaram a serem escritos em sequências de números hexadecimais. Em seguida o surgimento linguagens \en{assembly}, os códigos então começaram a ser escritos usando uma série de comandos mnemônicos que represetam uma conjunto de instruções em código de máquina. Essas linguagens de programação não são apropriadas para o desenvolvimento de aplicação em larga escala, portanto, o surgimento das linguagens de alto nível como C. Estas definem um conjunto de mecanismos que servem como abstrações das instruções de máquina (\en{loops}, condicionais, funções, etc) \cite{Bailey2005IntroductionC}.

  Construir estruturas de dados e algoritimos requer uma linguagem de programação de alto nível \cite{GoodrichTamassiaMount2011DataStructuresCpp}. Logo, para o desenvolvimento da aplicação móvel e dos \en{middlewares} de sincronização foram, portanto, escolhidas as linguagens de programação \en{Kotlin} e \en{Golang}, respectivamente.  



\subsubsection{LINGUAGEM DE PROGRAMAÇÃO: \en{KOTLIN}}

Para o desenvolvimento da aplicação móvel, sua construção foi realizada com o \en{software} \en{Android Studio} usando a linguagem de programação \en{Kotlin}. No desenvolvimento de uma aplicação é importante a escolha de uma plataforma, a plataforma escolhida para a aplicação foi o \en{Android}. Este é um sistema operacional para dispositivos móveis, sendo usado em \en{tablets}, celulares, dentre outros. \en{Android Studio} é o ambiente de desenvolvimento oficial de aplicações para a \en{Android}, pode ser usado com  \en{Java} e \en{Kotlin}. \en{Kotlin} é uma linguagem de programação pragmática que combina os paradigmas de Orientação a Objetos e Programação Funcional com foco em interoperabilidade; a interoperabilidade se dá pela qualidade da linguagem de usar códigos que foram escritos em \en{Java} em códigos escritos em \en{Kotlin} \cite{Resende2020KotlinAndroid}. 

Orientação a Objetos é uma paradigma de programação que determina a organização do código em objetos. Um objeto é gerado de uma classe, essa é a especificação de atributos e métodos que o objeto contém em sua definição; cada classe apresenta uma interface consistente para o restante do código sem entregar detalhes inapropriados \cite{GoodrichTamassiaMount2011DataStructuresCpp}. \en{Kotlin} como uma linguagem orientada a objetos adere a esse paradigma.


\subsection{LINGUAGEM DE PROGRAMAÇÃO: JAVA}

\en{Java} é uma linguagem de programação orientada a objetos, presente em diversos espaços. Os autores do \en{Java} organizaram os objetivos da linguagem em volta dos princípios: simplicidade, orientação a objetos, distribuição, robustez, segurança, neutralidade arquitetural, portabilidade, linguagem interpretada, performance, \en{multithreaded} e dinâmicidade. Como linguagem de programação, \en{Java} é frequentemente comparado com \en{C++}, onde ele se destaca na comparação por ter: \en{garbage collector} (não é necessário excluir as variáveis manualmente), o gerenciamento de apontadores de endereços de memória automático (no \en{C++} é algo que precisa ser levado em consideração no processo de construção código) e sua sintaxe mais amigável \cite{corejavafundamentals}.

No ambiente de sistemas distribuídos com \en{Java} pode-se destacar, como exemplo, o \en{CORBA} (\en{Common Object Request Broker Architecture}), esta tecnologia possibilita que sistemas totalmente distintos consigam operar em conjunto; \en{CORBA} complementa o \en{Java} por oferecer um ambiente de objetos distribuídos, a comunicação entre eles ocorre através de interfaces de programação \cite{oraclecorbajava}. Desse modo, é notável a versatilidade do \en{Java} como linguagem de programação; isso aliado à sua compatibilidade com o \en{Kotlin}, faz com que a escolha seja precisa para o contexto de desenvolvimento dos \en{middlewares} visados por esse texto. 



\subsection{MODELO CLIENTE-SERVIDOR}

A escolha dos componentes de um \en{software}, as interações entre eles e onde atuam, constitui a arquitetura de sistema. Na arquitetura de sistemas categoriza-se, basicamente, os sistemas em centralizados e descentralizados, essas categorizações se baseiam na natureza da relação cliente-servidor; Quando discutindo arquitetura de sistemas o modelo cliente-servidor é essencial. Desse modo, esse modelo define clientes que requisitam serviços de servidores, onde em sistemas centralizados é definido por: uma máquina cliente que implementa a parte do sistema voltada para o usuário; servidores contém o resto da aplicação (processamento e a camada de dados) \cite{tanenbaum2007dist}.

Em sistemas descentralizados essa organização difere, onde em sistemas centralizados a distribuição é alcançada, principalmente, por separar componentes logicamente em diferentes máquinas (distribuição vertical), nos descentralizados a distribuição é alcançada horizontalmente, isso é, clientes ou servidores podem ser divididos em partes equivalentes, cada uma responsável por processar sua parte do trabalho. Uma classe de arquitetura de sistema que implementa distribuição horizontal é \en{peer-to-peer} \cite{tanenbaum2007dist}.

\subsubsection{REDES \en{PEER-TO-PEER}}

Redes \en{peer-to-peer} são redes onde cada \en{node} atua como cliente e servidor. Essas redes podem ser de arquitetura estruturada e não estruturada. Nas estruturadas a rede é construída usando um procedimento determinista; um exemplo seria manter uma tabela com identificadores para todos os \en{nodes}, este mesmo identificador atua para categorizar os dados, assim quando um item é acessado na rede, é retornado o identificador do \en{node} responsável para ele. Em uma rede não estruturada os \en{nodes} conhecem uns aos outros de maneira quase aleatória, isso é, cada um mantém uma lista de vizinhos construída de maneira não determinista; a transmissão de dados ocorre da mesma forma \cite{tanenbaum2007dist}.



\subsection{SISTEMAS \en{OFFLINE-FIRST}}

\en{Offline-first} é um paradigma de programação que especifica, dentre outros aspectos, o comportamento de aplicações em relação a falta de conectividade \cite{OfflineFirst}. Um bom exemplo desse paradigma são os \en{Progressive Web Apps (PWAs)}, para uma aplicação ser um \en{PWA} esta deve englobar conceitos que formulam o modo que foi construída; dentre eles destaca-se, para esse trabalho, a independência de conexão \cite{BjornHansen2018PWAUnifiedMobile}. A arquitetura define, também, o armazenamento local de dados como componente essencial para sua implementação, e com ele a sincronização e a resolução de conflitos. \en{Offline-first} assume a existência da necessidade de conexão com alguma rede, sendo essa efêmera, e especifica como aplicações devem funcionar nesse contexto \cite{Guda2025OfflineFirstAndroidChat}, um exemplo deste tipo de conexão são as redes \en{peer-to-peer}.  

A oferta variável de recursos da computação móvel sugere que aplicações para esse contexto devem ser adaptáveis. Historicamente, adaptação é a troca de um recurso por outro ou pela qualidade do serviço prestado, isso é: técnicas de redução, filtragem e transformação de dados para trafegarem na rede; o processo de adaptabilidade é disparado automaticamente após a leitura do contexto. Contexto é toda informação relevante para aplicação que define seu comportamento; desse modo, as informações de contexto podem ser categorizadas em: físicas, relativas ao \en{hardware}; lógicas, relativas ao \en{software}; e, sociais, relativas à comunidade de usuários e seu contexto \cite{AugustinYaminGeyerBarbosa2001RequisitosAplicacoesMoveisDistribuidas}. Essas considerações são essenciais para garantir o funcionamento de aplicações que, ao seguirem o paradigma \en{Offline-first}, continuem funcionando mesmo na ausência de conexão, requisito essencial para esses tipos de soluções. 


\begin{comment}

Adicionar:
  - adicionar imagem do teorama CAP
  - adcionar imagem do resultado de impossibilidade FLP

Melhorar:
  - Fluidez entre paragráfos

Estrutura do conteúdo:
  - Explicar a relação entre comunicação assíncrona e sistemas Offline-First
  - Ressaltar o cenário ideal onde seria possível manter consistência e disponibilidade com tolerância a falhas
  - Mostrar como um sistema completamente assíncrono atende a essa demanda
  - Discorrer sobre o Teorema CAP para entender como não é possível manter consistência e disponibilidade com tolerância a falhas
  - Discorrer sobre como não é possível ter um sistema completamente assíncrono de acordo com o FLP
  - Concluir como essas ideias justificam a existência de estratégias de sincronização de dados

\end{comment}

\subsubsection{TEOREMA \en{CAP}}

Teorema \en{CAP} define um \en{tradeoff} entre consistência e disponibilidade em sistemas distribuídos não confiáveis. Apresentado por Eric Brewer através de um exemplo generalista de serviço \en{web} para demonstrar essa troca entre disponibilidade, consistência e tolerância a partição. Consistência sendo a capacidade do serviço de retornar a resposta correta para uma requisição; Disponibilidade é a garantia de toda requisição receber uma resposta, mesmo que eventualmente; e, tolerância a partição, este é a capacidade serviço de permanecer funcionando mesmo no caso de perda de comunicação \cite{CAPtheorem2012}.

No contexto de restrições de sistemas distribuídos, pode se destacar, ainda, o resultado de impossibilidade \en{FLP}; este afirma que em um sistema completamente assíncrono, isso é, em um sistema onde não há garantia de quando uma menssagem será entregue, não existe algoritmo determinístico capaz de encontrar consenso desde que o sistema assuma pelo menos uma falha de execução. Desse modo, entendendo o contexto com conectividade intermitente, é necessário, na construção de sistemas reais, relaxar a assincronicidade \cite{FLP1985}, por exemplo, ao aplicar mecanismos como: consistência eventual forte \cite{Gomes2017VerifyingSEC} ou \en{CRDTs}.

Assim, na construção da aplicação casos onde ocorrem desconexão de par devem ser resolvidos com aplicação de estratégias para manter os dados estruturados e coesos.


\subsubsection{ESTRATÉGIA DE SINCRONIZAÇÃO: \en{CRDT}}

Sincronização de dados em sistemas é uma tarefa essencial no estabelecimento de uma equivalência entre duas ou mais coleções de dados. Do mesmo modo, Tanenbaum \cite{tanenbaum2007dist} afirma que o principal desafio na replicação de dados (isso é, manter diferentes cópias de dados em servidores ou localidades diferentes) é assegurar a consistência das réplicas. 

Estratégias de sincronização estão presentes tanto em arquiteturas centralizadas e descentralizadas. Inicialmente, há os \en{CRDTs} (tipos de dado livre de conflito): algoritmos distribuídos que computam o estado de uma aplicação colaborativa do seu histórico de operações. Esse estado deve ser o mesmo para todos os usuários e ser computável sem a necessidade de ter um componente que centraliza as informações; um servidor central, por exemplo. Possui dois tipos distintos dependendo no que se baseiam: operações ou estados \cite{weidner2023crdtsurvey}.

\en{CRDTs} é um sistema replicado, e como tal precisa propagar atualizações para todas as replicas, há dois caminhos para alcançar esse objetivo: replicação baseada em estado e replicação baseada em operação. No primeiro, a estrutura de dados faz uso de uma função de união que integra o estado da replica remota na local; na segunda forma de replicação, a convergência dos estados ocorre através da propagação de operações entre replicas, portanto uma \en{CRDT} dessa categoria deve gerar, para cada operação, um gerador e um efetor; o primeiro responsável por gerar o segundo ao fim de uma operação; este segundo é, então, executado nas outras replicas para convergir os estados \cite{PreguicaBaqueroShapiro2018CRDT}.

\en{CRDTs} podem ser entendidas como resultado da evolução, até o momento, de estratégias de sincronização em sistemas distribuídos, a aplicação usando o \en{CRDT} de operação é o mais apropriado ao sistema de \en{Offline-First} produzido devido a capacidade das alterações em anotações serem traduzidas em operações (inserir e deletar).

Uma outra forma de sincronização que também faz uso do compartilhamento de estado é a transformação operacional (\en{OT}). Esta é voltada para arquiteturas centralizadas e é muito usada em editores em grupo; ela resolve problemas similares a \en{CRDTs}: consistência na manutenção de documentos compartilhados com tempo de resposta curto e edição paralela em ambientes distribuídos, por exemplo \cite{SunEllis1998OperationalTransformation}. Essa estratégia de sincronização é pesquisada desde o fim dos anos 80 e sua importância é inegável, desde artigos, videos, criação de jogos, comunicação \en{online} em geral. Os desenvolvimentos teóricos e provas experimentais de Ellis et al. que a originaram; atualmente, junto a \en{CRDT} é um dos métodos mais bem aceitos para resolver o problema de consistência de dados em arquiteturas distribuídas \cite{GadeaIonescuIonescu2020ControlLoopOT}.


\subsection{ESTRATÉGIA DE SINCRONIZAÇÃO: VETORES DE VERSÃO}

Outra abordagem para sincronização de dados são os vetores de versão. Vetores de versão são estruturas lógicas que codificam atualizações em sequências de pares formados pelo local e o número de atualizações daquele local; assim, em vetores de versão, duas réplicas são inconsistentes se, e apenas se, os seus vetores de versão são incomparáveis, isso é, cada réplica foi sujeita a atualizações que a outra não sofreu \cite{parker1983detection}.  Preguiça et al. (2010) expande ao afirmar que vetores de versão sintetizam histórias de causalidade ao guardar o prefixo dos eventos gerados em qualquer réplica especifica \cite{preguica2010dotted}.

A construção histórica dos vetores de versão inicia com os relógios de Lamport, onde é definido causalidade: se um evento A ocorre antes de um evento B, então o contador de A deve ser menor que o contador de B, contador sendo a marca temporal lógica dos eventos no sistema \cite{lamport1978time}. Relógios de Lamport não podem inferir concorrência, assim, Mattern introduz relógios vetoriais. As posições dos relógios vetoriais são referentes a processos, o valor é referente ao número de eventos ocorridos em cada processo; desse modo, é possível verificar a casualidade e concorrência dos eventos \cite{mattern1989virtual}. 

Portanto, relógios de Lamport codificam causalidade, relógios vetoriais detectam concorrência e vetores de versão adaptam o comportamento de relógios vetoriais para sistemas distribuídos. Ou seja, ao sintetizar o histórico causal das atualizações que ocorrem em uma réplica, vetores de versão explicitam quais são os estados correntes de cada uma; isso possibilita distinguir situações de equivalência, obsolescência e inconsistência entre réplicas \cite{almeida2002versionstamps}. 


\subsubsection{CONSISTÊNCIA EVENTUAL FORTE (\en{SEC})}

Consistência eventual forte define a consistência de réplicas sob o mesmo conjunto de atualizações, mesmo que essas atualizações sejam aplicadas em ordens distintas, desde que a casualidade seja respeitada. Interpretação de operações é, portanto, a base para a compreenção semâtica das operações entre as réplicas, pois garante coerência de contexto causal entre as réplicas. \en{SEC} assume que a interpretação de operações é correta, isso implica que as semâticas de domínio (especificações de cada sistema) devem ser preservadas para garantir a convergência \cite{Gomes2017VerifyingSEC}.

Gomes et al. (2017) apresentam exemplos de como a convergência de estado pode ser alcaçada por meio de regras que preservam a semâtica das operações em cenários concorrentes. Um desses exemplos é entitulado \en{Observable-Remove Set} (\en{OR Set}), este define remoções de elementos observáveis, isso é: cada elemento pode ser caracterizado por um identificador, a operação de remoção atua apenas sobre elementos nos quais o identificador é conhecido no momento de remoção; caso esses elementos removidos sejam adicionados novamente por uma operação concorrente, essa adição não será eliminada devido a discrepância de identificadores \cite{Gomes2017VerifyingSEC}. Esta implementação demonstra como a sincronização pode ser alcançada conforme as semânticas de domínio de cada sistema.

Ainda nos exemplos de Gomes et al. (2017), é descrito mecânismo de edição de texto que permite manter a caracteres excluídos, assim réplicas podem usar qualquer caractere como referência para inserir antes ou depois sem precisar considerar ações concorrentes. Isso pode ser aplicado na edição de anotações da aplicação móvel.

Nesse contexto, \en{CRDTs} garantem consistência eventual forte por especificar comutação de operações concorrentes, assim as réplicas podem convergir sem a necessidade de coordenação entre elas \cite{Shapiro2015CRDT}. 


\subsection{BANCO DE DADOS: \en{ROOM}}

Banco de dados são usados em diversos contextos (comercial, científico, pessoal) devido a necessidade do armazenamento persistente de dados; banco de dados são poderosos devido aos seus SGBDs (sistemas de gerenciamento de banco de dados) que provê as funcionalidades de persistência de dados, interface de modificação com uma liguagem de \en{query} (SQL, por exemplo), estas são linguagens de modificação de dados e da estrutura do banco de dados; e, genrenciamento de transação \cite{GarciaMolinaUllmanWidom2002DatabaseSystems}. Na aplicação móvel, o banco local será onde os dados sincronizados residiram.

\en{Room} é uma biblioteca de persistência que funciona sobre \en{SQLite}, fornecendo abstrações para permitir um acesso mais robusto ao banco de dados local \cite{AndroidJetpackRoom}. \en{SQLite} é um dos maiores sistemas de gerenciamento de banco de dados, estando presente em celulares, computadores, navegadores, televisores, etc; o motivo de seu sucesso pode ser creditado a: sua portabilidade para multiplas plataformas, ele é compacto e auto-contido, sua confiabilidade e velocidade \cite{GaffneyPrammerBrasfieldHippKennedyPatel2022}. \en{Room}, portanto, é a interface de banco de dados usada na aplicação móvel. 


\subsubsection{AMBIENTE DE TESTES: \en{DOCKER}}

\en{Docker} é uma plataforma aberta para desenvolvimento, publicação e execução de aplicações. Ele oferta recursos para empacotar e executar uma aplicação em ambiente com isolamento fraco, isso é, um contêiner. \en{Docker} também oferece ferramentas para orquestração de contêineres, onde é possível manipular a criação, destruição desses; também é possível isolar ambientes, simulando redes distintas. \en{Docker} é usado a partir da manipulação de seus objetos: contêineres, \en{networks}, volumes, imagens, etc. Cada objeto atua sobre um sistema operacional definido no contêiner, por exemplo, \en{Debian}, \en{Ubuntu}, \en{Fedora}, dentre outros \cite{DockerOverviewDocs}.



\subsection{TRABALHOS RELACIONADOS}

Os trabalhos relacionados a essa pesquisa estão presentes no contexto de sincronização de dados em sistemas distribuídos. Desse modo, pode-se destacar o trabalho de Lamport, responsável por formalizar a noção de causalidade entre eventos em sistemas computacionais, que se torna a base para rastrear dependências entre eventos distribuídos \cite{lamport1978time}. Mattern oferece uma progressão natural ao introduzir a noção de tempo virtual e relógios vetoriais, que possibilita a detecção de concorrência entre eventos \cite{mattern1989virtual}. Assim, Parker et al. exibe uma perspectiva central para sistemas \en{Offline-First} onde afirma: quando réplicas evoluem de forma independente enquanto desconectadas, conflitos são o resultado natural desse modo de operar, não uma exceção; então sugerem um mecanismo de rastreamento de versões que permite comparar dados e distinguir se duas cópias são equivalentes \cite{parker1983detection}.

Nessa linha de manter as réplicas consistentes há o trabalho de Demers et al.; nele é investigado algoritmos epidêmicos, isso é: algoritmos que propagam dados de modo similar à propagação de epidemias. No trabalho é discutido métodos como \en{direct mail}, \en{rumor mongering} e \en{anti-entropy} \cite{demers1987epidemic}; algoritmos epidêmicos são uma boa escolha para redes \en{peer-to-peer} devido ao seu caráter descentralizado.

Para exemplos práticos nota-se os trabalhos que aplicam o paradigma \en{Offline-First}. Um desses trabalhos é Requisitos para o Projeto de Aplicações Móveis Distribuídas \cite{AugustinYaminGeyerBarbosa2001RequisitosAplicacoesMoveisDistribuidas} que apresenta o modelo ISAM para construir um sistema de produção de aplicativos para o contexto móvel; publicado em 2001 permanece relevante até os dias atuais devido o seu trabalho na categorização de contexto. Outro trabalho é \en{Implementation of Offline-First Architectures for Android Internet-Based Chat Systems} \cite{Guda2025OfflineFirstAndroidChat} publicado em 2025 segue a mesma linha, propõe uma estrutura para a criação de um \en{chat} para redes locais, não lida com redes descentralizadas mas oferece comentários úteis na sincronização de mensagens. 

Desse modo, os trabalhos citados apontam uma trajetória conceitual: definição de conceitos básicos como causalidade, segue para mecanismos de detecção de incosistência e propagação de dados de forma descentralizada; finalmente, alcança aplicações práticas de sistemas projetados para funcionar no contexto de dispositivos móveis. É nesse percurso que este trabalho se insere, como um estudo das estratégias de sincronização de dados em redes \en{peer-to-peer} com aplicações que adotam o paradigma \en{Offline-First}. 
